%%=============================================================================
%% Methodologie
%%=============================================================================

\chapter{\IfLanguageName{dutch}{Praktische Uitwerking}{}}%
\label{ch:methodologie}

\section{\IfLanguageName{dutch}{Printplaten}{}}%
\label{sec:PCB}

De kabeltester is opgebouwd uit een omhulsel en drie printplaten: \textit{Controller board}, \textit{expansion board} en \textit{interconnection board} (Figuur \ref{fig:BDTester}). Deze zijn ontworpen door mijn interne mentor Johan Scheirs. 

\begin{figure}[h!]
    \centering
    \captionsetup{width=0.8\textwidth}
    \includegraphics[width=0.8\textwidth]{BDTester.png}
    \caption[Blokdiagram kabeltester]{Blokschema van de kabeltester}
    \label{fig:BDTester} 
\end{figure}

\subsection{\IfLanguageName{dutch}{Controller board}{}} 
\label{ssec:controllerboard}

Dit is de centrale printplaat van de kabeltester (Figuur \ref{fig:ControllerBoard}). Hierop bevinden zich de ESP32, level shifter, EEPROM, teach-knop, connectors voor voeding, \textit{expansion board} en case met leds en knop (Figuur \ref{fig:BDController}). De gebruikte ESP32 in dit prototype is de ESP32-S3-DevKitM-1U van Espressif \autocite{Espressif2026} en is het brein achter de kabeltester. Deze zal op een toekomstige release vervangen worden door de chip zelf met een kleinere \textit{footprint}. De signalen dat deze microcontroller stuurt naar het \textit{expansion board}, passeren eerst de level shifter om de uitgaande signalen om te zetten van 3,3 V naar 5 V aangezien de ESP en alle schuifregisters aan de 5 V rail hangen. De EEPROM die gebruikt wordt is de 24LC64. Het wordt gestuurd via de microcontroller met I2C en kan tot acht kilobytes opslaan. 

\begin{figure}[h!]
    \centering
    \captionsetup{width=0.8\textwidth}
    \includegraphics[width=0.8\textwidth]{ControllerBoard.jpg}
    \caption[Controller board]{Bovenaanzicht van het controller board met microcontroller gemonteerd en draadjes naar voorkant omhulsel}
    \label{fig:ControllerBoard} 
\end{figure}

\begin{figure}[h!]
    \centering
    \captionsetup{width=0.8\textwidth}
    \includegraphics[width=0.8\textwidth]{BDController.png}
    \caption[Blokdiagram controller board]{Blokschema van het controller board}
    \label{fig:BDController} 
\end{figure}


\subsection{\IfLanguageName{dutch}{Expansion board}{}} 
\label{ssec:expansionboard}

De tweede printplaat bevat de schuifregisters (Figuur \ref{fig:ExpansionrBoard}). Per bord vind je vier paar bestaande uit één SIPO en één PISO schuifregisters bij elkaar. Hierbij is de eerste parallel out bit verbonden met de eerste parallel in bit. Deze verbinding takt ook af naar een de eerste pin in een \textit{8-pin female header}. Daarnaast heeft dit bord ook twee connectors voor \textit{flatcables} waarvan, als je van boven af kijkt, de linkse verbonden wordt met het controller bord en de rechtse optioneel kan verbonden worden met een volgende \textit{expansion board}. Zo kan je de lengte van de \textit{daisy chain} uitbreiden in het geval dat de kabel meer dan 32 pins nodig heeft.

\begin{figure}[h!]
    \centering
    \captionsetup{width=0.8\textwidth}
    \includegraphics[width=0.8\textwidth]{ExpansionBoard.jpg}
    \caption[Expansion board]{Opgelost expansion board}
    \label{fig:ExpansionrBoard} 
\end{figure}

De schuifregisters die gekozen zijn voor dit project zijn de 74HC4094 voor het uitsturen en de HEF4021B voor het inlezen. Deze worden gestuurd via SPI met een gezamenlijke klok. Daarnaast heeft de 74HC4094 een \textit{output enable} (OE) en een \textit{strobe} pin. De eerste dient voor de parallelle uitgangen door te laten of tegen te houden. De \textit{strobe} kan gebruikt worden om tijdens het schuiven de nieuwe schuivende waarde rechtstreeks op de uitgangen te zetten. Voor de HEF4021B te sturen hebben we ook nog een extra signaal nodig voor de \textit{Parallel / Serial Control} (P/S) of \textit{Parallel Load} (PL) te sturen, benaming is anders afhankelijk van producent. Omdat het schuifregister niet tegelijk serieel kan schuiven en de parallelle pinnen uit kan lezen, moeten we kort wisselen tussen beide zodat de waarden kunnen uitlezen.

\begin{figure}[h!]
    \centering
    \captionsetup{width=0.8\textwidth}
    \includegraphics[width=0.8\textwidth]{BDExpansion.png}
    \caption[Blokdiagram expansion board]{Blokschema van het expansion board}
    \label{fig:BDExpansion} 
\end{figure}

\begin{figure}[h!]
\centering
\captionsetup{width=0.8\textwidth}
\includegraphics[width=0.4\textwidth]{PCBFout_U2.png}
\includegraphics[width=0.4\textwidth]{PCBFout_U4.png}
\includegraphics[width=0.4\textwidth]{PCBFout_U6.png}
\includegraphics[width=0.4\textwidth]{PCBFout_U8.png}
\caption[Design fouten expansion board]{Design fouten PISO schuifregisters expansion board}
\label{fig:ExpansionFout} 
\end{figure}

In dit prototype zijn ook nog enkele defecten gevonden (Figuur \ref{fig:ExpansionFout}). Bij de PISO schuifregisters zijn de massa pinnen (VSS) niet verbonden met het grondvlak en zweven ze nog. Daarnaast is de \textit{serial out} (Q8) van schuifregister U6 verbonden aan de twee volgende U4 en U2. Dit zorgt ervoor dat de \textit{serial out} van U4 open blijft terwijl deze met de \textit{serial in} van U2 moest verbonden zijn. Deze problemen zijn dan met enkele jumper draadjes en wat soldeerwerk opgelost op de geleverde printplaatjes.

\subsection{\IfLanguageName{dutch}{Interconnection board}{}} 
\label{ssec:interconnectionboard}

Als laatste hebben we het interconnection board (Figuur \ref{fig:IOBoard}). Deze heeft dezelfde dimensies als het expansion board en dient om de pinnen die verbonden zijn aan de ingangen en uitgangen van de schuifregisters te verbinden met bruikbare sockets waar je de verschillende connectors van de kabel kan insteken (Figuur \ref{fig:BDIO}). Dit betekent dat er enkel een nieuw \textit{interconnection board} moet worden ontworpen als de kabeltester gebruikt wordt voor een nieuwe kabel. Op het prototype is het \textit{interconnection board} gemaakt op maat van een kabel van Trince. Ook deze moest bij aankomst nog in orde gebracht worden omdat er twee pads nog overlapten met de trace van hun buren door het veranderen van ronde pads naar vierkante (Figuur \ref{fig:IOFout}).

\begin{figure}[h!]
    \centering
    \captionsetup{width=0.8\textwidth}
    \includegraphics[width=0.8\textwidth]{IOBoard.jpg}
    \caption[IO board]{Interconnection board voor een specifieke kabel van Trince.}
    \label{fig:IOBoard} 
\end{figure}

\begin{figure}[h!]
    \centering
    \captionsetup{width=0.8\textwidth}
    \includegraphics[width=0.8\textwidth]{BDIO.png}
    \caption[Blokdiagram expansion board]{Blokschema van een mogelijk interconnection board}
    \label{fig:BDIO} 
\end{figure}

\begin{figure}[h!]
    \centering
    \captionsetup{width=0.8\textwidth}
    \includegraphics[width=0.4\textwidth]{PCBFout_XYHOME.png}
    \includegraphics[width=0.4\textwidth]{PCBFout_XAXIS.png}
    \caption[Kortsluitingen IO board]{Kortsluitingen bij connectors interconnectie board}
    \label{fig:IOFout} 
\end{figure}

\newpage
\section{\IfLanguageName{dutch}{Standalone functionaliteit}{}}%
\label{sec:CTStandalone}

Van de kabeltester wordt verwacht dat deze autonoom kan werken. Hiervoor moet je enkel een kabel hebben die je wil testen en een kabeltester die een \textit{golden sample} opgeslagen heeft om mee te vergelijken.

\subsection{\IfLanguageName{dutch}{Analyse}{}}%
\label{sec:AnalyseStandalone}

De standalone functionaliteit van de kabeltester heeft twee functies: het testen van een kabel door een operator (Use Case \ref{uc:Test}) en het aanleren/opslaan van een \textit{golden sample} door een beheerder aan de hand van de 'Teach' (Use Case \ref{uc:Teach}). Het testen is heel voor de hand liggend. Een operator hoeft enkel de kabel te verbinden met de juiste connectoren op het \textit{I/O board} en vervolgens op de knop op de behuizing duwen. Hierna zal de kabeltester de verbonden kabel inlezen en vergelijkt deze met de referentie die hij heeft opgeslagen. Indien de patronen overeenkomen zal de groene led opblinken anders zal de rode led oplichten (Figuur \ref{fig:StandaloneSequentie}). Het aanleren van de referentie vraagt echter om de behuizing open te schroeven. Dit komt omdat de knop om een kabel aan te leren op het \textit{controller board} staat. Hiervoor is gekozen zodat een werknemer niet zomaar de \textit{golden sample} zou overschrijven met een andere, mogelijks foute, kabel. Net zoals de test zal de \textit{teach} functie de kabel eerst inlezen en deze dan opslaan in de EEPROM. 

\begin{figure}[h!]
    \centering
    \captionsetup{width=0.8\textwidth}
    \begin{tcolorbox}[colback=white, colframe=odisee-navy, title=\textbf{Use Case: Kabel testen}, width=\textwidth, sharp corners,breakable]
        \textbf{Kenmerken:}
        \begin{itemize}[noitemsep, label={}]
            \item \textbf{Primaire actor:} De operator
            \item \textbf{Preconditie:} De kabeltester moet zijn ingeschakeld.
            \item \textbf{Postconditie:} De kabeltester geeft aan of de kabel correct is.
        \end{itemize}
        
        \textbf{Hoofdscenario:}
        \begin{enumerate}[noitemsep]
            \item De operator wil zijn kabel testen.
            \item De operator sluit de kabel aan op het expansion board. (\textit{DR-CABLEPINS})
            \item De operator drukt op de externe drukknop om de test te starten.
            \item De kabeltester controleert de kabelstructuur.
            \item De kabeltester toont het resultaat via de leds.
        \end{enumerate}
        
        \tcblower % Dit trekt een lijntje tussen de bovenkant en de domeinregels
        \textbf{Domeinregel DR-CABLEPINS:} Maximaal 256 pinnen, 32 schuifregisters, 8 uitbreidingsborden.
    \end{tcolorbox}
    \caption[Use Case Test]{Use Case: Kabel testen}
    \label{uc:Test} 
\end{figure}

% --- Use Case 2: Nieuwe kabel inleren ---
\begin{figure}[h!]
    \centering
    \captionsetup{width=0.8\textwidth}
    \begin{tcolorbox}[colback=white, colframe=odisee-navy, title=\textbf{Use Case: Een nieuwe kabel inleren}, width=\textwidth, sharp corners, label={uc:Teach}]
        \textbf{Kenmerken:}
        \begin{itemize}[noitemsep, label={}]
            \item \textbf{Primaire actor:} Een beheerder (administrator)
            \item \textbf{Preconditie:} De behuizing van de kabeltester is geopend en de kabeltester is ingeschakeld.
            \item \textbf{Postconditie:} De verbindingen van de nieuwe kabel zijn in kaart gebracht en opgeslagen.
        \end{itemize}
        
        \textbf{Hoofdscenario:}
        \begin{enumerate}[noitemsep]
            \item De beheerder wil een nieuwe kabel inleren.
            \item De beheerder sluit de kabel aan op het interconnectiebord. (\textit{DR-CABLEPINS})
            \item De beheerder drukt op de interne drukknop ``TEACH''.
            \item De kabeltester brengt verbindingen in kaart.
            \item De kabeltester slaat de configuratie op in het geheugen.
            \item De PASS-led kleurt groen bij succes.
        \end{enumerate}
        
        \textbf{Alternatieve Scenario's:}
        \begin{itemize}[noitemsep]
            \item \textbf{5A. Geheugenfout:} FAIL-led kleurt rood; use case eindigt zonder postconditie.
        \end{itemize}
        
        \tcblower 
        \textbf{Domeinregel DR-CABLEPINS:} Maximaal 256 pinnen, 32 schuifregisters, 8 uitbreidingsborden.
    \end{tcolorbox}
    \caption[Use Case Teach]{Use Case: Een nieuwe kabel inleren}
    \label{uc:Teach} 
\end{figure}

\begin{figure}[h!]
    \centering
    \captionsetup{width=0.8\textwidth}
    \includegraphics[width=0.8\textwidth]{SequenceDiagramV1.png}
    \caption[Sequentiediagram standalone]{Sequentiediagram van standalone functionaliteit}
    \label{fig:StandaloneSequentie} 
\end{figure}

\newpage
\subsection{\IfLanguageName{dutch}{Uitwerking}{}}%
\label{ssec:Uitwerking}

\subsubsection{\IfLanguageName{dutch}{Architectuur}{}}%
\label{sssec:Architectuur}

De kabeltester is opgebouwd uit verschillende componenten (Figuur \ref{fig:Blokdiagram}). Deze zijn opgedeeld in de applicatie laag, drivers en \textit{board support packages} (BSP). De bovenste laag, de applicatielaag, bevat de kernlogica van het systeem. Hierin worden de testroutines uitgevoerd, meetresultaten geanalyseerd en beslissingen genomen over de status van de aangesloten kabel. Om dit te bereiken zijn er drivers geschreven die de brug vormen tussen de applicatie en de fysieke hardware. Dit gebeurt door gebruik te maken van de in ESP32 meegeleverde BSPs die \textit{low-level} acces bieden aan de hardware. De console, API, menu, microrl en UART spelen bij de autonome versie geen rol. Hun functionaliteit  wordt besproken in hoofdstuk \ref{sssec:AanpassigenKabeltester}.

\begin{figure}[h!]
    \centering
    \captionsetup{width=0.8\textwidth}
    \includegraphics[width=0.8\textwidth]{Blokdiagram.png}
    \caption[Blokdiagram]{Overzicht van architectuur zowel standalone als MES-versie}
    \label{fig:Blokdiagram} 
\end{figure}

\paragraph{\IfLanguageName{dutch}{Applicatie laag}{}}

\begin{itemize}[]
    \item \textbf{Main:} Deze functie aangeboden in ESP-IDF wordt bij opstart aangeroepen en start de application.
    \item \textbf{App:} Dit is de primaire taak binnen de applicatie. Als eerste initialiseert het de \textit{CableOps} en de taken \textit{StatusApp} en \textit{Console}. Daarna wacht het tot een nieuw commando in de wachtrij terecht komt. Wanneer dit gebeurt roept de app de desbetreffende functie op.
    \item \textbf{CableOps:} In dit component staat de logica om een kabel in te lezen, te vergelijken en op te slaan.
    \item \textbf{StatusApp:} Deze taak behandelt de leds in de behuizing van de kabeltester. Door net als de App een wachtrij in de gaten te houden.
\end{itemize}

\paragraph{\IfLanguageName{dutch}{Drivers}{}}

\begin{itemize}[]
    \item \textbf{Shifter:} Aan de hand van de SPI package behandelt deze driver het inschuiven en uitlezen van de schuifregisters; 74HC4094 en HEF4021B. Om dit te doen moet het wel eerst via GPIO de output pinnen van de level shifter en de 74HC4094 hoog zetten. Daarna kan het nieuwe patroon ingeschoven en het vorige uitgelezen worden. Tussen elke transactie moet de driver ook voor een korte tijd de P/S hoog zetten om de nieuwe signalen in het register op te slaan.
    \item \textbf{DebugLed:} In deze driver wordt de ingebouwde led op de microcontroller aangesproken. Op de ESP32-S3-DevKitM-1U kan deze een RGB waarde krijgen. Deze driver is voor de werking niet noodzakelijk en kan weggelaten worden als er overgegaan wordt naar een printplaat met een ESP32 rechtstreeks op gemonteerd.
    \item \textbf{I2CEeprom:} Met behulp van de meegeleverde I2C package regelt deze driver het contact met de EEPROM. Lezen van de opgeslagen data is altijd mogelijk maar voor schrijven moet eerst de \textit{Write Protect} (WP) pin laag gezet worden. Als dat niet het geval is stuurt de EEPROM een \textit{Negative Acknowledgement} of NACK terug.
    \item \textbf{Button:} Deze driver leest de GPIO pinnen die overeenkomen met de interne en externe knop aan de hand van interrupts. Deze voegt dan het overeenkomende commando toe aan de wachtrij.
    \item \textbf{StatusLeds:} Via de ledc stuurt deze driver een pulse naar de PASS-led of de FAIL-led afhankelijk van de parameter.
\end{itemize}

Hieronder verdiepen we ons wat dieper op de belangrijke bestanden.

\paragraph{\IfLanguageName{dutch}{CableOps}{}}%

De belangrijkste functionaliteiten bevinden zich in dit component. Hier worden de \textit{test} en \textit{teach} functie gedefinieerd die in de App opgeroepen worden. Beide scannen eerst de kabel en slaan die op in een globale variabele '\textit{g\_cable[CABLE\_CURRENT]}'. Vervolgens zal de test deze waarde vergelijken met de referentie die zich bevind in de variabele '\textit{g\_cable[CABLE\_REFERENCE]}'. Deze referentie wordt bij opstarten opgehaald uit de EEPROM. De \textit{teach} zal in plaats van te vergelijken deze waarde overschrijven met het nieuwe patroon waarna het deze ook in de EEPROM zal opslaan.

Om een kabel te scannen kijken we eerst hoe veel \textit{expansion boards} er aan de kabeltester hangen. Op basis van die informatie wordt de lengte van het bitpatroon bepaald. In de huidige versie, omdat maar één board werkt, staat deze waarde \textit{NUM\_REGISTERS hardcoded} op vier. We creëren enkele variabelen aan om te gebruiken als buffers en resetten de huidige kabel door de waarden op nul te zetten (Codefragment \ref{cf:scan_cable_1}).

\begin{listing}[H]
    \begin{minted}[fontsize=\footnotesize, breaklines, baselinestretch=0.95]{C}
        /**
        * @brief Scans the cable by iterating through all possible bit combinations and
        *        recording unique connection patterns.
        * @return ESP_OK if the scan completes successfully, otherwise an error code.
        */
        static esp_err_t scan_cable() {
            static uint8_t tx_buffer[NUM_REGISTERS];
            static uint8_t current_rx[NUM_REGISTERS];
            static uint8_t rx_buffer[NUM_REGISTERS];
            g_cable[CABLE_CURRENT].count = 0;
            g_cable[CABLE_CURRENT].size = 0;
            
            memset(g_cable[CABLE_CURRENT].pattern, 0, sizeof(g_cable[CABLE_CURRENT].pattern));
            memset(rx_buffer, 0, NUM_REGISTERS);
    \end{minted}
    \caption[scan\_cable pt 1]{Het initialiseren van buffers}
    \label{cf:scan_cable_1}
\end{listing}

Vervolgens starten we de scan. Hiervoor moeten we de uitgangen van de \textit{level shifter} eerst activeren. Dit gebeurt via de functie \lstinline|shifter\_enable\_levelshifter()| die de gpio-pin verbonden met de OE hoog maakt. Vervolgens starten we een lus die in totaal \lstinline|TOTAL\_BITS + 1| keer uitvoeren. Hierbij is \textit{TOTAL\_BITS} gelijk aan 32 omdat we vier schuifregisters hebben met ruimte voor één woord (8 bits). We doen één extra iteratie omdat het resultaat van een ader hoog te zetten telkens één iteratie achter loopt. Dit komt omdat onze schuifregisters de SPI poort delen. In elke iteratie wordt de tx\_buffer gevuld met een bitpatroon waar maar één bit hoog staat. Elke keer dat de teller van de lus incrementeert zal de hoge bit één positie verschuiven. Dan sturen we het patroon door naar de schuifregisters en slaan terwijl het resultaat op in \textit{current\_rx} (Coderfragment \ref{cf:scan_cable_2}).

\begin{listing}[H]
    \begin{minted}[fontsize=\footnotesize, breaklines, baselinestretch=0.95]{C}
            ESP_LOGI("SCAN", "Starting scan...");
            
            shifter_enable_levelshifter();
            
            // We run TOTAL_BITS + 1 to catch the "trailing" result of the final bit
            for (int i = 0; i <= TOTAL_BITS; i++) {
                // Set 1 bit high for every iteration
                memset(tx_buffer, 0, NUM_REGISTERS);
                if (i < TOTAL_BITS) { // Only set a bit if we are within the actual
                    // hardware range
                    tx_buffer[i / 8] = (1 << (i % 8));
                }
                
                shifter_update_registers(tx_buffer, current_rx, false);
    \end{minted}
    \caption[scan\_cable pt 2]{Code van de functie die de kabels inleest deel 2.}
    \label{cf:scan_cable_2}
\end{listing}

Omdat de het resultaat van de eerste transactie niet van belang is verwerken we deze niet en slaan we het over. Bij de volgende iteraties starten we met de lijst om te keren. Dit is nodig omdat de PISO schuifregisters van achter naar voor verbonden zijn waardoor de bytes in het patroon gespiegeld zijn rond de middelste. Daarna controleren we of het resultaat verschillend is van de \textit{tx\_buffer} en het uniek is. Dit doen we om plaats te besparen want het opslaan van een open connectie of gelijke patronen zou geen meerwaarde geven. Na de lus zetten we de \textit{level shifter} weer uit (Codefragment \ref{cf:scan_cable_3}). 

\begin{listing}[H]
    \begin{minted}[fontsize=\footnotesize, breaklines, baselinestretch=0.95]{C}
                if (i > 0) {
                    // Reverse the order of bits to match the physical layout of the
                    // pins
                    utils_reverse_array(current_rx, NUM_REGISTERS);
                    ESP_LOG_BUFFER_HEX("RX_ARRAY", current_rx, NUM_REGISTERS);
                    
                    // Check if the result is different from the input (indicating a
                    // connection)
                    if (memcmp(current_rx, rx_buffer, NUM_REGISTERS) != 0) {
                        
                        // The result is different from the input, but is it unique?
                        bool is_duplicate = false;
                        for (int k = 0; k < g_cable[CABLE_CURRENT].count; k++) {
                            if (memcmp(g_cable[CABLE_CURRENT].pattern[k], current_rx,
                            NUM_REGISTERS) == 0) {
                                is_duplicate = true;
                                break;
                            }
                        }
                        
                        // Save only if it's a new, unique connection pattern
                        if (!is_duplicate && g_cable[CABLE_CURRENT].count < TOTAL_BITS) {
                            memcpy(g_cable[CABLE_CURRENT].pattern[g_cable[CABLE_CURRENT].count],
                            current_rx, NUM_REGISTERS);
                            ESP_LOGI("SCAN", "Found Connection! Pattern saved at index %d",
                            g_cable[CABLE_CURRENT].count);
                            g_cable[CABLE_CURRENT].count++;
                        }
                    }
                }
                ESP_LOGI("SCAN", "Bit %d", i);
                ESP_LOG_BUFFER_HEX("TX_ARRAY", tx_buffer, NUM_REGISTERS);
                
                memcpy(rx_buffer, tx_buffer, NUM_REGISTERS);
                utils_delay_ms(10);
            }
            ESP_LOGI("SCAN", "Scan Complete.");
            // Disable TXB0108 level shifter (active high)
            shifter_disable_levelshifter();
            g_cable[CABLE_CURRENT].size = g_cable[CABLE_CURRENT].count * NUM_REGISTERS;
            return ESP_OK;
        }
    \end{minted}
    \caption[scan\_cable pt 3]{Code van de functie die de kabels inleest deel 3.}
    \label{cf:scan_cable_3}
\end{listing}

\subsection{\IfLanguageName{dutch}{Problemen}{}} 

Niet alles verliep zoals verwacht. Bij het aansluiten van een tweede \textit{expansion board} achter het eerste is er op te merken dat het resulterende bitpatroon niet overeenkomt met wat er verwacht wordt. Hetzelfde fenomeen is op te merken wanneer we een te lange \textit{flatcable} tussen het  \textit{controller} en \textit{expansion board}. Er zijn verschillende hypotheses wat de oorzaak van deze problemen kunnen zijn:

\begin{itemize}[]
    \item De \textit{level shifter} beschikt mogelijk over onvoldoende \textit{drive strength} om meerdere ingangen correct aan te sturen, waardoor de uitgangsspanning onder de logische drempelwaarde van de schuifregisters kan zakken.
    \item Bij een langere kabel kan de verbindingslijn zich meer als een transmissielijn gedragen, wat leidt tot terugkaatsende signalen en ongewenste oscillaties op de flanken.
    \item Er is een te hoge capacitieve belasting. Naarmate de afstand dat de signalen afleggen en voor elke ingang waaraan deze verbonden is, wordt parasitaire capaciteit toegevoegd. Dit kan de flanken van de signalen in steilheid verminderen.
\end{itemize}

Het uitbreiden met meerdere \textit{expansion boards} is een belangrijke eis van de kabeltester. Daarom moet er in de toekomst verder onderzocht worden wat de exacte oorzaak is. 

\newpage
\section{\IfLanguageName{dutch}{MES functionaliteit}{}}%
\label{sec:CTMES}

\subsection{\IfLanguageName{dutch}{Analyse}{}}%
\label{ssec:AnalyseMES}

Voor de kabeltester te gebruiken in combinatie met de MES moeten we een micro-USB aansluiten op de voorziene UART-poort van het Devkit en de USB-A kant bij de MES. Hierna kan de operator van de MES de \textit{workinstruction} starten en door de stappen lopen. Wanneer hij bij de stap komt van het test-script zal de MES via de seriële connectie de commando's \textit{test}, \textit{get-ref} en \textit{get-cur} sturen. In volgorde zullen zij de test uitvoeren, de kabel referentie ophalen en de huidige kabel ophalen. Hierna zal met deze informatie een YAML-bestand gemaakt worden. Daarmee kan een kabel diagram gecreëerd worden met behulp van Wireviz. Waarna deze beide aan de operator getoond worden (Use Case \ref{uc:MES}).

% --- Use Case 1: Kabel Testen ---
\begin{figure}[h!]
    \centering
    \captionsetup{width=0.8\textwidth}
    \begin{tcolorbox}[colback=white, colframe=odisee-navy, title=\textbf{Use Case: Een kabel testen via MES}, sharp corners, label={uc:MES}]
        \textbf{Kenmerken:}
        \begin{itemize}[noitemsep, label={}]
            \item \textbf{Primaire actor:} De operator
            \item \textbf{Ondersteunende actoren:} Kabeltester en QMS
            \item \textbf{Preconditie:} Kabeltester en MES verbonden via seriële communicatie; workinstruction geopend.
            \item \textbf{Postconditie:} Resultaat opgeslagen in QMS.
        \end{itemize}
        
        \textbf{Hoofdscenario:}
        \begin{enumerate}[noitemsep]
            \item Operator werkt voorgaande stappen werkinstructie af.
            \item Operator sluit kabel aan op expansion board. (\textit{DR-CABLEPINS})
            \item Operator start stap met python script.
            \item Systeem stuurt commando over seriële lijn.
            \item Kabeltester controleert structuur en stuurt PASS/FAIL.
            \item Systeem haalt meetgegevens op en vergelijkt met referentie.
            \item Systeem genereert en toont Wireviz-diagrammen.
            \item Operator vergelijkt diagrammen en systeem laat stap slagen.
            \item Operator werkt werkinstructie af; resultaat opgeslagen in QMS.
        \end{enumerate}
        
        \textbf{Alternatieve Scenario's:}
        \begin{itemize}[noitemsep]
            \item \textbf{3A. Seriële connectie faalt:} Systeem toont error en laat stap falen en gaat naar stap 9.
            \item \textbf{8A. Kabel incorrect:} Systeem laat stap falen en gaat verder naar 9.
        \end{itemize}
        
        \tcblower 
        \textbf{Domeinregel DR-CABLEPINS:} Maximaal 256 pinnen, 32 schuifregisters, 8 uitbreidingsborden.
    \end{tcolorbox}
    \caption[Use Case MES]{Use Case: Een kabel testen via de MES}
    \label{uc:MES} 
\end{figure}

\newpage
\subsection{\IfLanguageName{dutch}{Uitwerking}{}}%
\label{ssec:UitwerkingMES}

\subsubsection{\IfLanguageName{dutch}{Aanpassingen kabeltester}{}}%
\label{sssec:AanpassigenKabeltester}

Om de kabeltester te laten communiceren over seriële communicatie werden de console, API, menu en microrl componenten toegevoegd (Figuur \ref{fig:Blokdiagram}). Met deze samen kun je via seriële communicatie de kabeltester sturen, exact wat de MES nodig heeft om te werken. De kabeltester stuurt de resultaten steeds terug als JSON-object.

\begin{itemize}[]
    \item \textbf{Console:} In dit component bevind zich de \textit{FreeRTOS} taak die de seriële poort leest en ernaar schrijft.
    \item \textbf{API:} Dit component voegt de commando's toe aan het menu en dient als een brug tussen de menu en de App.
    \item \textbf{Menu:} Deze driver bied methodes aan om de menu te vullen met mogelijke commando's.
    \item \textbf{Microrl\autocite{microrl}:} Deze externe driver bied hulp bij het creëren van de \textit{command line interface} (CLI).
\end{itemize}

Hieronder verdiepen we ons wat dieper op de belangrijke bestanden.

\paragraph{Console}

De primaire taak van de console is het inlezen van de seriële communicatie. Het resultaat hiervan stuurt het naar het menu. Daarnaast maakt de console ook de JSON-objecten. Er zijn twee types: één voor het resultaat van de test en \textit{teach} en een ander voor een kabel op te halen. Het resultaat van een test of \textit{teach} ziet er als volgt uit (Codefragment \ref{cf:TestJSON}). Hier zal de waarde bij de sleutel `res` gelijk zijn aan `PASS` of `FAIL`.

\begin{listing}[H]
    \begin{minted}[fontsize=\footnotesize, breaklines, baselinestretch=0.95]{text}
        {
            "res":"PASS"
        }
    \end{minted}
    \caption[Test JSON]{Voorbeeld JSON van test of \textit{teach}.}
    \label{cf:TestJSON}
\end{listing}

Het resultaat van de commando's om de referentie of huidige kabel op te halen bestaan uit drie sleutels en waarden (Codefragment \ref{cf:CableJSON}): de data die het patroon als \textit{base64} gecodeerd heeft staan, de \textit{size} dat bestaat uit het aantal connecties in het patroon en de \textit{checksum} van de data om te controleren dat de data tijdens het versturen niet \textit{corrupt} is geraakt.

\begin{listing}[H]
    \begin{minted}[fontsize=\footnotesize, breaklines, baselinestretch=0.95]{text}
        {
            "data":"AQAAgAJJABAEAAgAEAAEACAAEACAACAAAJAAAgAgAQAAAIAI",
            "size":9,
            "crc":1548549446
        }
    \end{minted}
    \caption[Kabel JSON]{Voorbeeld JSON van het ophalen van een kabelpatroon.}
    \label{cf:CableJSON}
\end{listing}


\paragraph{Menu}

Met behulp van \textit{Microrl} kan er een eenvoudige CLI gemaakt worden. Als je via seriële communicatie naar de kabeltester stuurt kan je commando's uitvoeren die aan het menu zijn toegevoegd. Deze kunnen rechtstreeks toegevoegd worden aan het hoofdmenu of er kan een submenu gemaakt worden voor verschillende functionaliteiten. Het menu bied zelf ook een hulp commando aan die meer info geeft over de mogelijke commando's binnen het menu waarin je je begeeft. Voor eenvoud zijn alle commando's die de MES nodig heeft, toegevoegd aan het hoofdmenu (Codefragment \ref{cf:help}).

\begin{listing}[H]
    \begin{minted}[fontsize=\footnotesize, breaklines, baselinestretch=0.95]{text}
        CT> help
        I (9431) MENU: Available commands:
        I (9431) MENU: console     : console functions
        I (9431) MENU: test        : Test a cable
        I (9431) MENU: teach       : Learn a new cable
        I (9431) MENU: get_cur      : Get current pattern
        I (9431) MENU: get_ref      : Get reference pattern
        CT> 
    \end{minted}
    \caption[Menu - help commando]{Resultaat van het help commando van het menu.}
    \label{cf:help}
\end{listing}


\subsubsection{\IfLanguageName{dutch}{Workinstruction}{}}%
\label{sssec:WI}

Een \textit{workinstruction} bestaat uit een Jira ticket waar een stappenplan is aan verbonden (Te zien in de rode kader n°1 op figuur \ref{fig:WIPage}). Optioneel kan daar een \textit{package} aan toegevoegd worden die python-scripts bevatten (Groene kader n°2). Indien een package is toegevoegd, kan vervolgens een script aangeroepen worden in een stap. Dit kan gebruikt worden om automatisatie toe te voegen aan de \textit{workinstruction}. 
In het voorbeeld zijn wordt zowel in stap 1 als stap 2 een script opgeroepen met 

`\lstinline|<qity:script name="cable-tester-test" args="/dev/ttyUSB0" interactive="true"/>|`

waar de naam bepaald welk script het is. Voor \textit{`cable-tester-test`} wordt ook het argument '/dev/ttyUSB0' meegegeven wat aangeeft op welke USB poort de kabeltester is aangesloten. Deze huidige versie gaat er wel van uit dat de kabeltester het enige USB toestel aangesloten op de MES is. 

\begin{figure}[h!]
    \centering
    \captionsetup{width=0.8\textwidth}
    \includegraphics[width=0.8\textwidth]{WIPage.png}
    \caption[Workinstruction pagina]{Jira pagina van een simpele kabeltester \textit{workinstruction}}
    \label{fig:WIPage} 
\end{figure}

Van elke \textit{workinstruction} kunnen er meerdere instanties uitgevoerd worden (Figuur \ref{fig:DiagramWI}). Wanneer een instantie afgewerkt is, kunnen de resultaten in de tab \textit{test runs} gevonden worden. Deze bevatten de informatie welke stap misliep en de bestanden die opgeslagen zijn zoals de diagrammen. Zo kan er later teruggekeken worden wat het probleem was.

\begin{figure}[h!]
    \centering
    \captionsetup{width=0.8\textwidth}
    \includegraphics[width=0.8\textwidth]{DiagramWI.png}
    \caption[Diagram workinstruction]{Diagram over de werking van een \textit{workinstruction}}
    \label{fig:DiagramWI} 
\end{figure}

Een \textit{workinstruction} ziet er als volgt uit op de MES (Figuur \ref{fig:MESInterface}). Bovenaan heb je de naam van het Jira ticket. Onder de titel zie je de naam van het MES toestel onder 'Equipment ID', de operator onder de header met dezelfde naam, de vooruitgang van de stappen onder status en de unit ID. Wanneer je een \textit{workinstruction} wil afwerken is het verplicht om een unit ID mee te geven. Bij elke stap staat er een \textit{FAIL} en \textit{PASS} knop. Hiermee kan de operator laten weten of een stap slaagde of niet. Bij stappen gelinkt aan een script zal deze keuze bij het eindigen van de logica door het script zelf gebeuren. 

\begin{figure}[h!]
    \centering
    \captionsetup{width=0.8\textwidth}
    \includegraphics[width=0.8\textwidth]{MESInterface.png}
    \caption[MES interface]{Interface van MES bij \textit{workinstruction} 'Test cable'.}
    \label{fig:MESInterface} 
\end{figure}


\subsubsection{\IfLanguageName{dutch}{Python package}{}}%
\label{sssec:PythonPackage}

De Python \textit{package} bestaat uit verschillende folders. Hiervan zijn sommige gelijk voor alle \textit{workinstructions}. Deze gemeenschappelijke folders bevatten scripts voor testen, debuggen en mes functionaliteiten. Daarnaast is er ook een folder voor mijn scripts genaamd \textit{ctclient}. Hierin bevinden zich de volgende Python bestanden:

\begin{itemize}[]
    \item \textbf{cli:} Dit is het eerste bestand dat opgeroepen wordt wanneer het script \textit{cable-tester-test} wordt uitgevoerd. Hierin worden de argumenten die aan het script zijn toegevoegd verwerkt. Daarna zal het de test of \textit{teach} uitgevoerd worden afhankelijk van wat er van de \textit{workinstructions} verwacht wordt.
    \item \textbf{cable\_tester\_client:} Hierin bevindt zich de klasse \textit{CableTesterClient}. Deze bevat de methodes test en \textit{teach} die elk via de interne methode \textit{\_send} de nodige commando's naar de kabeltester en geeft de resultaten door aan \textit{wireviz\_generator}
    \item \textbf{wireviz\_generator:} Om de patronen om te zetten naar een diagram gebruiken we Wireviz. Deze verwacht een YAML bestand met de data die het moet omvormen. Dit script creëert zo een bestand en maakt er vervolgens grafische bestanden van die het opslaat in de \textit{out} folder van de MES. 
    \item \textbf{cable\_tester\_app:} Hoewel de MES al een grafische interface heeft is het momenteel niet mogelijk om images uit de \textit{package} weer te geven in de stappen. Daarom wordt er via de Python \textit{library Qt} een nieuw venster gemaakt die de gegenereerde diagrammen weergeeft. 
\end{itemize}

Hieronder verdiepen we ons wat dieper op de belangrijke bestanden.

\paragraph{\IfLanguageName{dutch}{Cable\_tester\_client}{}}%

In dit bestand staan de definities van de test en \textit{teach} functies van de \textit{client}. Om een test uit te voeren worden de commando's test, \textit{get\_ref} en \textit{get\_cur} opeenvolgend uitgestuurd met behulp van de interne methode \textit{\_send} (Codefragment \ref{cf:_send}). Als resultaat krijgt de test dan een \textit{dictionary} terug met de data. Hierbij is echter het patroon nog als \textit{base64} gecodeerd en moet dus gedecodeerd worden. Wanneer dit in orde is maken we een nieuwe \textit{dictionary} aan die we meegeven aan de generator. Vanaf dat de diagrammen opgesteld en opgeslagen zijn, opent de test vervolgens een nieuw venster die beide diagrammen naast elkaar toont (Codefragment \ref{cf:test_ctclient}). Het enige verschil tussen beide is dat de \textit{teach} maar één patroon ophaalt en dus ook maar één diagram genereert en weergeeft. 

\begin{listing}[H]
    \begin{minted}[fontsize=\footnotesize, breaklines, baselinestretch=0.95]{Python}
        def test(self) -> int:
            res = self._send(b"test\r\n")
            
            res_cur = self._send(b"get_cur\r\n")
            current_pattern = self.decodeBase64(res_cur)
            
            res_ref = self._send(b"get_ref\r\n")
            reference_pattern = self.decodeBase64(res_ref)
            
            result = res.get("res", "unknown")
            
            test_dict: TestResponse = {
                "result": result,
                "current_pattern": current_pattern,
                "current_count": len(current_pattern),
                "reference_pattern": reference_pattern,
                "reference_count": len(reference_pattern),
            }
            
            # print(f"Test Dictionary: {test_dict}")
            generate_cable_diagram(test_dict, self._out_dir)
            
            self._app.display_diagram(
            f"{self._out_dir}/cable_correct.svg", f"{self._out_dir}/cable_fault.svg"
            )
            ret_code = self._app.exec()
            match ret_code:
                case MESReturnCode.NO_ERROR:
                    pass
                case MESReturnCode.CANCELED:
                    print("RFID programming step was canceled")
                case _:
                    print(
                    f"Cable tester client application ended with exit code {ret_code}"
                    )
            
            ret: int = 0 if result == "PASS" else 1
            return ret
    \end{minted}
    \caption[Ctclient test]{Code van de test methode.}
    \label{cf:test_ctclient}
\end{listing}

De \textit{\_send} methode stuurt het commando door en leest wat er via de seriële connectie terug wordt gestuurd. Aangezien het resultaat een JSON-object bevat zoekt de functie naar de accolades die het begin en einde van het object afbakenen. Wanneer dit is gevonden wordt het omgevormd tot een \textit{dictionary} en teruggegeven als resultaat. Dit allemaal staat in een \textit{try} blok die mogelijke errors zal opvangen (Codefragment \ref{cf:_send}).

\begin{listing}[H]
    \begin{minted}[fontsize=\footnotesize, breaklines, baselinestretch=0.95]{Python}
        def _send(self, data: bytes) -> Any:
            port = self._config["port"]
            raw: str | None = None
            extracted_data: Any = None
            try:
                with serial.Serial(**self._config) as ser:
                    ser.write(data)
                    ser.flush()
                    buffer = ""
                    while True:
                        # Read available bytes (more efficient than 1 by 1)
                        response = ser.read(ser.in_waiting or 1)
                        if not response:
                            break
                    
                        # Decode and append to our buffer
                        buffer += response.decode("utf-8", errors="ignore")
                        
                        # Look for something that looks like a JSON object: { ... }
                        # The dot (.) doesn't match newlines by default, so we use re.DOTALL
                        match = re.search(r"(\{.*\})", buffer, re.DOTALL)
                        
                        if match:
                            json_string = match.group(1)
                            # Validate that it is actually valid JSON
                            extracted_data = json.loads(json_string)
                            # print("Successfully extracted JSON:")
                            # print(extracted_data)
                            break
            
            except serial.serialutil.SerialException as exc:
                raise CableTesterException(f"Serial port error on '{port}': {exc}") from exc
            except PermissionError as exc:
                raise CableTesterException(
                f"Permission denied opening '{port}': {exc}"
                ) from exc
            except FileNotFoundError as exc:
                raise CableTesterException(f"Port not found '{port}': {exc}") from exc
            except json.JSONDecodeError:
                print(f"FAILED TO DECODE: '{raw}'")
                raise CableTesterException(f"Invalid JSON response from device on '{port}'")
            
            return extracted_data
    \end{minted}
    \caption[Ctclient \_send]{Code van de interne send methode.}
    \label{cf:_send}
\end{listing}

\paragraph{\IfLanguageName{dutch}{Wireviz\_generator}{}}%

Om een diagram te maken heeft \textit{Wireviz} een YAML-bestand nodig. Dit bestand is opgedeeld in connectors (Xx), \textit{cables} (Wx) en \textit{connections}. Onze kabeltester kan echter de fysieke kabels en connectors op het \textit{interconnection board} niet achterhalen, enkel de connecties krijgen we. Daarom wordt er aan de stap al een YAML-bestand als bijlage toegevoegd die de connectors van de kabel al bevat. Dit bestand wordt dus ideaal gecreëerd tegelijkertijd met het ontwerp van het \textit{interconnection board}. 

Nadat dit bestand opgehaald en omgevormd is naar een \textit{dictionary} worden de connectors hun pinnen toegewezen aan hun locatie in het bitpatroon. Vervolgens kijken we hoe de \textit{wires} zijn opgebouwd door de patronen te analyseren. Zo komen we voor elke pin te weten aan hoeveel andere pinnen ze hangen. Hierna wordt er voor zowel de referentie als de geteste kabel een compleet YAML-bestand gemaakt die wordt opgeslagen in de \textit{out} folder en vervolgens met \textit{Wireviz} wordt gerenderd. (Codefragment \ref{cf:generate_cable_diagram})

\begin{listing}[H]
    \begin{minted}[fontsize=\footnotesize, breaklines, baselinestretch=0.95]{Python}
        def generate_cable_diagram(
        response: TestResponse, out_dir: str, io_yml: str | None = None
        ) -> None:
            io_path = resolve_io_yml(out_dir, io_yml)
            io_connectors = load_io_definition(io_path)
            pm = build_pin_map(io_connectors)
            
            ref_wires = parse_pattern(response["reference_pattern"])
            cur_wires = parse_pattern(response["current_pattern"])
            
            # Reference diagram — all green
            correct_path = os.path.join(out_dir, CORRECT_YML)
            save_yaml(
                build_wireviz_yaml(
                    ref_wires,
                    title="Reference — all connections correct",
                    io_connectors=io_connectors,
                    pin_map=pm,
                ),
                correct_path,
            )
            _render(correct_path)
            
            # Fault diagram — colour-coded diff
            diff = diff_patterns(ref_wires, cur_wires)
            fault_path = os.path.join(out_dir, FAULT_YML)
            save_yaml(
                build_wireviz_yaml(
                    diff,
                    title="Fault analysis (GN=ok  OG=missing  RD=short)",
                    io_connectors=io_connectors,
                    pin_map=pm,
                ),
                fault_path,
            )
            _render(fault_path)
            
            _print_summary(ref_wires, cur_wires, diff, out_dir, io_connectors, pm)
    \end{minted}
    \caption[Wireviz generate cable diagram]{Code van de functie die de diagrammen maakt.}
    \label{cf:generate_cable_diagram}
\end{listing}

De diagrammen zien er als volgt uit (Figuur \ref{fig:cable_diagrammen}). Omdat we niet kunnen weten uit hoeveel deelkabels de effectieve kabel bestaat gaan we in de diagrammen elke unieke connectie een eigen \textit{wire} (Wx) geven. Hierbij zal soms een \textit{wire stub} (WxS) gemaakt in combinatie met een \textit{splice connector} (SPx) omdat meerdere aders niet kunnen vertrekken van één pin in \textit{Wireviz}. Aangezien de referentie de juiste connecties bevat, zullen alle aders in het diagram groen gekleurd zijn. Bij de geteste kabel is ideaal ook alles groen wat zou betekenen dat het in orde is. Indien er een fout in zit zullen oranje aders connecties tonen die we verwachten maar niet aanwezig zijn en rode aders tonen kortsluitingen die we niet willen zien. 

\begin{figure}[h!]
    \centering
    \captionsetup{width=0.8\textwidth}
    \includegraphics[width=0.8\textwidth]{cable_correct.png}
    \includegraphics[width=0.8\textwidth]{cable_fault.png}
    \caption[Kabel diagrammen]{Voorbeeld van de gegenereerde kabel diagrammen gebaseerd op de voorbeeld kabel van Trince.}
    \label{fig:cable_diagrammen} 
\end{figure}

Momenteel kan de \textit{splice connector} wel een vertekend beeld geven. Omdat bijvoorbeeld het linkse deel van SP1 een fout bevat, is het verwarrend dat het rechtse deel wel groen aangeeft. In een toekomstige update zou in plaats van groen deze zijde in een grijze kleur weergegeven worden. Op die manier toont het wel aan dat het deels in orde is maar dat de linkerzijde invloed heeft op de rechter.

